% Direction C: "Merged Alloy + clause matrix". Stage 1 and Stage 1.5 share one blob, because they share
% the solver and three of four clauses: a clause MATRIX (rows = clauses, columns = stage) shows exactly what
% changes, and the Stage-1 graphs loop back in as the pin. Passes = one horizontal chevron timeline with
% labels alternating above/below. gem5 = a chip package: the gray die is the defense build, the PINS are the
% non-invasive harness (it touches only the boundary). Finding = the bug under a magnifier lens.
\documentclass[tikz,border=0pt]{standalone}
\usepackage{newtxtext,newtxmath}
\usepackage{graphicx}
% Shared shape vocabulary for the big-figure design mocks (BIGFIG_SHAPE_DESIGN.md).
% Every tikzpicture that uses these must set [x=1pt,y=1pt]: all numbers below are pt.
\usetikzlibrary{calc,fit,shapes.symbols,shapes.misc,shapes.geometric,decorations.pathmorphing,arrows.meta}
\def\defshade{black!12}   % the one gray: defense-specific
\tikzset{
  >=latex,
  every node/.style={font=\footnotesize, inner sep=0pt,
     execute at begin node={\hyphenpenalty=10000\exhyphenpenalty=10000\relax}},
  q/.style={font=\scriptsize\itshape},             % qualifier: multiplicity, mode, arrow label
  sm/.style={font=\scriptsize},                    % small roman label
  blob/.style={draw, rounded corners=7pt, line width=0.6pt},          % stage container
  btitle/.style={fill=white, inner xsep=2pt, inner ysep=1pt, anchor=west}, % title sitting on the blob border
  slab/.style={draw, rounded corners=2.5pt, minimum height=10.5pt, inner xsep=3pt, align=left, anchor=north west},
  ghost/.style={slab, dash pattern=on 1.4pt off 1.1pt},               % "same as Stage 1"
  pill/.style={draw, rounded corners=4.5pt, minimum height=10.5pt, inner xsep=3.5pt, align=center},  % formula pill
  tool/.style={draw, rounded rectangle, minimum height=17pt, inner xsep=5pt, align=center, line width=0.6pt},  % solver / engine
  chev/.style={draw, signal, signal to=east, signal from=west, signal pointer angle=125,
               minimum height=19pt, align=center, inner xsep=1pt, font=\scriptsize},
  chevL/.style={draw, signal, signal to=west, signal from=east, signal pointer angle=125,
               minimum height=19pt, align=center, inner xsep=1pt, font=\scriptsize},
  chevD/.style={draw, signal, signal to=south, signal from=north, signal pointer angle=140,
               minimum width=15pt, minimum height=10pt, align=center, inner sep=0pt, font=\scriptsize\bfseries},
  dspec/.style={fill=\defshade},
  dot/.style={circle, fill, inner sep=0.75pt},
  station/.style={circle, draw, fill=white, minimum size=8.5pt, inner sep=0pt, font=\scriptsize\bfseries},
  box/.style={draw, rounded corners=3pt, align=center, inner xsep=3pt, inner ysep=2.2pt},
}

% ---------- glyphs drawn inside a card centered at the scope origin (card 11 x 14) ----------
\newcommand{\gGraph}{% execution graph: 3 instruction nodes in spo + one rf arc (matches fig a)
  \draw[line width=0.3pt] (-1.8,3.6) -- (-1.8,-3.6);
  \fill (-1.8,3.6) circle (0.95) (-1.8,0) circle (0.95) (-1.8,-3.6) circle (0.95);
  \draw[line width=0.35pt, -{Latex[length=1.5pt,width=1.5pt]}] (-1.0,3.8) .. controls (3.6,3.2) and (3.6,-3.2) .. (-0.9,-3.5);}
\newcommand{\gAsm}{% program text
  \foreach \y/\l in {4/5.2,2/3.8,0/5.2,-2/2.8,-4/4.4} \draw[line width=0.35pt] (-3,\y) -- ++(\l,0);}
\newcommand{\gAnn}{\node[font=\tiny] at (0,0.2) {\{\,\}};}
\newcommand{\gTrace}{% pipeline trace: staggered stage bars (O3PipeView-like)
  \foreach \y/\a/\b in {4/-3.6/0.4,1.6/-2.4/1.6,-0.8/-1.2/2.8,-3.2/0/3.6} \draw[line width=0.9pt] (\a,\y) -- (\b,\y);}

% ---------- card stack: #1 name, #2 (x,y) of the front card's center, #3 number of cards, #4 glyph macro ----------
% back cards sit up-right of the front card; anchors #1-n (top of the whole stack), #1-s, #1-e, #1-w
\def\cw{11}\def\ch{14}\def\co{1.8}
\newcommand{\cardstack}[4]{%
  \begin{scope}[shift={#2}]
    \foreach \k in {#3,...,1} {\pgfmathsetmacro\o{\co*(\k-1)}
      \draw[fill=white, rounded corners=0.9pt] (-\cw/2+\o,-\ch/2+\o) rectangle ++(\cw,\ch);}
    #4
    \pgfmathsetmacro\top{\ch/2+\co*(#3-1)}\pgfmathsetmacro\rgt{\cw/2+\co*(#3-1)}
    \path (\co*#3/2-\co/2,\top) coordinate (#1-n) (0,-\ch/2) coordinate (#1-s)
          (\rgt,\co*#3/2-\co/2) coordinate (#1-e) (-\cw/2,0) coordinate (#1-w);
  \end{scope}}
% stack of program + annotation PAIRS: each card is two half-cards clipped together
\newcommand{\pairstack}[3]{% #1 name, #2 (x,y), #3 n
  \begin{scope}[shift={#2}]
    \foreach \k in {#3,...,1} {\pgfmathsetmacro\o{\co*(\k-1)}
      \draw[fill=white, rounded corners=0.9pt] (-11+\o,-7+\o) rectangle ++(10.4,14);
      \draw[fill=white, rounded corners=0.9pt] (0.6+\o,-7+\o) rectangle ++(10.4,14);
      \draw[line width=0.5pt] (-1.4+\o,4.5+\o) -- (1.4+\o,4.5+\o);}  % clip joining the two halves
    \begin{scope}[shift={(-5.8,0)}]\gAsm\end{scope}
    \begin{scope}[shift={(5.8,0)}]\gAnn\end{scope}
    \pgfmathsetmacro\top{7+\co*(#3-1)}
    \path (-5.8,7) coordinate (#1-asm-n) (5.8,7) coordinate (#1-ann-n) (-5.8,-7) coordinate (#1-asm-s) (5.8,-7) coordinate (#1-ann-s)
          (11,0) coordinate (#1-e) (-11,0) coordinate (#1-w) (0,\top) coordinate (#1-n) (0,-7) coordinate (#1-s);
  \end{scope}}

% ---------- defense-specific source of truth: gray folded document; #1 name, #2 (x,y), #3 glyph or empty ----------
\newcommand{\sot}[3]{%
  \begin{scope}[shift={#2}]
    \draw[fill=\defshade] (-5.5,-7) -- (5.5,-7) -- (5.5,4) -- (2.5,7) -- (-5.5,7) -- cycle;
    \draw (2.5,7) -- (2.5,4) -- (5.5,4);
    \if\relax\detokenize{#3}\relax
      \foreach \y in {1.8,-0.8,-3.4} \draw[line width=0.3pt] (-3.2,\y) -- (3.2,\y);
    \else \node[font=\tiny] at (0,-1.2) {#3};\fi
    \path (0,7) coordinate (#1-n) (0,-7) coordinate (#1-s) (5.5,0) coordinate (#1-e) (-5.5,0) coordinate (#1-w);
  \end{scope}}

% ---------- bug glyph (#1 position, #2 scale, #3 color) ----------
\newcommand{\buggl}[3]{%
  \begin{scope}[shift={#1}, scale=#2, line width=0.45pt, draw=#3, fill=#3]
    \foreach \y in {1.0,0,-1.0} {\draw (-2.6,\y-0.45) -- (2.6,\y+0.45); \draw (-2.6,\y+0.45) -- (2.6,\y-0.45);}
    \fill (0,-0.2) ellipse (1.35 and 1.9);
    \fill (0,2.15) circle (0.8);
    \draw (-0.35,2.7) -- (-1.2,3.7) (0.35,2.7) -- (1.2,3.7);
  \end{scope}}
% magnifier (#1 center of lens, #2 radius)
\newcommand{\magn}[2]{\draw[line width=0.6pt] #1 circle (#2); \draw[line width=1.1pt] ($#1+(-45:#2)$) -- ++(-45:#2*0.9);}
% legend swatch
\newcommand{\swatch}[1]{\draw[fill=\defshade] #1 rectangle ++(6,6);}

\begin{document}
\begin{tikzpicture}[x=1pt,y=1pt]
% ---------------- knobs ----------------
\def\LX{24}
\def\PL{56.3}
\def\PR{251.7}
\def\IL{62}
\def\SX{69}
\def\MW{112}    % matrix width
\def\ca{130}    % x of the Stage-1 column
\def\cb{160}    % x of the Stage-1.5 column
\useasboundingbox (0,1) rectangle (252,-344);
\coordinate (lxc) at (\LX,0); \coordinate (spx) at (214,0);
\tikzset{cell/.style={font=\scriptsize}, row/.style={slab, text width=\MW pt-6pt}}

% ---------------- legend ----------------
\swatch{(0,-20.6)}
\node[q, anchor=north west, align=left] at (7.5,-13) {defense-specific\\sources of truth};

% ================= Stages 1 + 1.5 (one blob) =================
\def\T{-6}
\node[q] at (\ca,\T-11) {Stage 1};
\node[q] at (\cb,\T-11) {Stage 1.5};
\node[row] (r1) at (\IL,\T-17) {\textbf{1}\enspace well-formed};
\node[row] (r2) at ($(r1.south west)+(0,-2)$) {\textbf{2}\enspace minimality};
\node[row, dspec] (r3) at ($(r2.south west)+(0,-2)$) {\textbf{3}\enspace leakage};
\node[row] (r4) at ($(r3.south west)+(0,-2)$) {\textbf{4}\enspace pin};
\foreach \r in {r1,r3} {\node[dot, inner sep=1pt] at (\r -| \ca,0) {}; \node[dot, inner sep=1pt] at (\r -| \cb,0) {};}
\node[cell] at (r2 -| \ca,0) {strict};
\node[cell] at (r2 -| \cb,0) {loosened};
\node[cell] at (r4 -| \ca,0) {--};
\begin{scope}[shift={(\cb,0)}, shift={(r4.center -| 0,0)}, scale=0.62]\gGraph\end{scope}
\draw[line width=0.6pt] ($(r1.north east)+(3,0)$) -- ($(r4.south east)+(3,0)$);
\coordinate (m1) at ($(r1.north east)!0.5!(r4.south east)$);
\node[tool] (A) at (m1 -| spx) {Alloy Analyzer\\[-1.5pt]{\scriptsize Kodkod + SAT}};
\draw[->] ($(r1.north east)!0.5!(r4.south east)+(3,0)$) -- (A.west);
% Stage-1 run -> exhaustive graphs -> loop back as the pin of each Stage-1.5 run
\cardstack{s1}{(200,\T-77)}{3}{\gGraph}
\draw[->] (s1-n |- A.south) -- (s1-n);
\node[q, anchor=east] at (199,\T-58) {1};
\draw[->] (s1-w) -- (\cb,\T-77) -- (r4.south -| \cb,0);
\node[q, anchor=north west, align=left] at (\IL+2,\T-71) {exhaustive Stage-1 graphs;\\each one is pinned};
\draw[blob] (\PL,\T) rectangle (\PR,\T-90);
\node[btitle] at (\IL+2,\T) {\textbf{Stages 1, 1.5}\enspace Abstract tests \textit{(Alloy)}};
\sot{ds}{(r3.west -| lxc)}{}
\node[sm, anchor=north, align=center] at ($(ds-s)+(0,-1.5)$) {Defense\\specification};
\draw[->] (ds-e) -- (r3.west);

% ---- output: Stage-1.5 graphs ----
\def\G{-102}
\draw (230,{\T-17-48/2-8.5}) -- (230,\G) -- (\SX,\G);
\node[q, anchor=west] at (232,\T-57) {1.5};
\cardstack{o2}{(\SX,\G-16)}{4}{\gGraph}
\draw[->] (\SX,\G) -- (o2-n);
\node[anchor=west, align=left] at (\SX+16,\G-14) {More flexible leaking execution graphs\\[-1pt]{\scriptsize\itshape multiple per Stage-1 graph; non-exhaustive;}\\[-2pt]{\scriptsize\itshape parallelizable SAT solving}};

% ================= Stage 2: one chevron timeline =================
\def\T{-136}
\draw[->] (o2-s) -- (\SX,\T-29);
\def\cy{\T-35.5}
\foreach \i/\x/\t in {1/75.5/{PC order,\\state},2/100.7/{pick\\instructions},3/125.8/{speculation\\primitives},4/151/{interleave\\chains},5/176.2/{assign regs +\\mem; reserve},6/201.3/{SSA + LLVM,\\PC labels},7/226.5/{asm +\\annotations}} {
  \node[chev, minimum width=27pt, minimum height=13pt, font=\scriptsize\bfseries] (p\i) at (\x,\cy) {\i};
  \pgfmathparse{int(mod(\i,2))}
  \ifnum\pgfmathresult=1
    \node[sm, anchor=north, align=center] at ($(p\i.south)+(0,-1.8)$) {\t};
  \else
    \node[sm, anchor=south, align=center] at ($(p\i.north)+(0,1.8)$) {\t};
  \fi
}
\draw[blob] (\PL,\T) rectangle (\PR,\T-62);
\node[btitle] at (\SX+6,\T) {\textbf{Stage 2}\enspace Concretization \textit{(Python, LLVM)}};

% ---- output 3 ----
\def\G{-206}
\draw (p7.east) -- (246,\cy) -- (246,\G) -- (\SX+9,\G);
\pairstack{o3}{(\SX+9,\G-17)}{4}
\draw[->] (\SX+9,\G) -- (\SX+9,\G-5.5);
\node[anchor=west, align=left] at (\SX+31,\G-15) {Programs (x86 asm) + annotations\\[-1pt]{\scriptsize\itshape one pair per branch mode $\times$ interleave variant}};

% ================= Stage 3: gem5 as a chip =================
\def\T{-243}
\def\J{-234.5}
% package (gem5) with pins (harness) and a gray die (defense build)
\draw[rounded corners=2.5pt, line width=0.6pt] (66,\T-16) rectangle (146,\T-60);
\node[anchor=north west] at (69,\T-18.5) {\includegraphics[height=6pt]{gem5logo.pdf}};
\node[draw, rounded corners=1.5pt, dspec, minimum width=60pt, minimum height=21pt, align=center, font=\scriptsize] (die) at (110,\T-42) {defense build\\under test};
\foreach \x in {84,96,108,120,132} {\draw[line width=0.6pt] (\x,\T-16) -- (\x,\T-12.5); \draw[line width=0.6pt] (\x,\T-60) -- (\x,\T-63.5);}
\foreach \y in {-30,-42,-54} \draw[line width=0.6pt] (146,\T+\y) -- (149.5,\T+\y);
\node[q, anchor=west] at (150,\T-21) {harness pins};
\draw[->] (o3-asm-s) -- (o3-asm-s |- 0,\T-16);
\draw (o3-ann-s) -- (o3-ann-s |- 0,\J) coordinate (jn);
\draw[->] (jn) -- (jn |- 0,\T-12.5);
\node[dot] at (jn) {};
\node[q, anchor=west] at ($(jn |- 0,\T-7)+(2.5,0)$) {directives};
% trace -> inspection -> bug under the lens
\cardstack{tr}{(160,\T-42)}{1}{\gTrace}
\draw[->] (149.5,\T-42) -- (tr-w);
\node[box, minimum height=22pt, anchor=west, font=\scriptsize] (I) at (171,\T-42) {Inspection\\program};
\draw[->] (tr-e) -- (I.west);
\draw[->] (jn) -- (I.north |- jn) -- (I.north);
\node[q, anchor=west] at ($(I.north |- 0,\T-7)+(1.5,0)$) {labels};
\draw[line width=0.9pt, fill=white] (230,\T-38) circle (9.5);
\draw[line width=2pt] ($(230,\T-38)+(-45:9.5)$) -- ++(-45:7);
\buggl{(230,\T-38.5)}{1.25}{black}
\node[font=\footnotesize\bfseries, anchor=north] at (226,\T-50.5) {Bug?};
\draw[->] (I.east) -- (220.5,\T-42);
% sources
% SimSpect goal: the implementation should satisfy the specification (fills the otherwise empty lane)
\draw[dashed, ->] (\LX,-79) -- (\LX,-276);
\node[q, rotate=90, anchor=south] at (\LX-2,{(-79+-276)/2}) {\textbf{SimSpect goal:} validate security guarantee in implementation};
\sot{gb}{(\LX,\T-42)}{\textless/\textgreater}
\node[sm, anchor=north, align=center] at ($(gb-s)+(0,-1.5)$) {Defense's\\gem5 build};
\draw[->] (gb-e) -- (die.west);
\sot{tm}{(\LX,\T-78)}{}
\node[sm, anchor=north, align=center] at ($(tm-s)+(0,-1.5)$) {Transmission\\mechanisms};
\draw (tm-e) -- (I.south |- tm-e);
\draw[->] (I.south |- tm-e) -- (I.south);
\draw[->] (108,\T-78) -- (108,\T-63.5);
\node[dot] at (108,\T-78) {};
\draw[blob] (\PL,\T) rectangle (\PR,\T-85);
\node[btitle] at (\SX+21,\T) {\textbf{Stage 3}\enspace Execution \textit{(gem5)}};
\end{tikzpicture}
\end{document}
